\cotefile{34}
\scanpage{84}
\mspage{1}
% archive sheet number: 83
% top heading
\uline{Exposé} III \uline{le Théorème de \unsure{monodromie} par voie géométrique}

1. \uline{Comparaison de la cohomologie des \unsure{deux|2} fibres} \deleted{\unsure{spéciales}}\inserted{\uline{géométriques}} \deleted{\unsure{et|de} la fibre} \uline{spéciale et générale} : la \uline{suite exacte des cycles évanescents}.

% upper-left diagram; the manuscript strokes have no clear arrowheads
\diagram{\begin{tikzcd}[ampersand replacement=\&,column sep=small,row sep=small]
X \arrow[d,no head] \& X_K \arrow[d,no head] \& X_{\bar K} \arrow[d,no head] \\
S \arrow[dr,no head] \& K \arrow[dr,no head] \& \bar K \\
\& \widetilde S \arrow[r,no head] \& \widetilde K
\end{tikzcd}}

Notations de I 1.1. On suppose donné $X$ \uline{propre} sur $S$, d'où $X_{\unsure{\eta|y}}=X_K$ sur $\eta=\operatorname{Spec}K$, et on se propose de \deleted{\unsure{définir|déterminer}} donner des \unsure{relations} entre $H^i(X_{\bar s})$ et $H^i(X_{\bar\eta})$ (coefficients $\Lambda_\nu=\mathbb Z/\ell^{\nu+1}\mathbb Z$ ou $\mathbb Z_\ell=\varprojlim_\nu\Lambda_\nu$). On verra que cette question est liée intimement à celle des opérations de $\unsure{\pi|\Pi}$ sur $H^i(X_{\bar\eta})$ ; [par exemple on \unsure{désirerait|disposerait} un \unsure{homomorphisme|isomorphisme} $H^i(X_{\bar s})\simeq H^i(X_{\bar\eta})^\pi$].

Quitte à remplacer $S$ par $\widetilde S$, on peut supposer pour simplifier les notations que $S$ est \uline{strictement} \deleted{\unsure{hensélien}} local (et \unsure{$\simeq$|même} complet \unsure{si on y tient} \unsure{\dots}, quand $\ell\ne\operatorname{car}K$). On a
\[
\bar K=\varinjlim_{\unsure{i}}\deleted{\unsure{K}}K_\alpha\deleted{\unsure{K_\alpha}}
\tag*{(1.1)}
\]
$K_\alpha$ les sous-extensions finies de $K$, d'où
\[
X_{\bar\eta}=\varprojlim_{\unsure{i}}X_\alpha\deleted{K_\alpha}\quad\text{où}\quad X_\alpha=X_{K_\alpha},\quad\text{et}
\tag*{(1.2)}
\]
\[
H^{\unsure{*|d}}(X_{\bar\eta},\Lambda_\nu)=\varinjlim_\alpha H^*(X_{\bar K_\alpha},\Lambda_\nu)
\tag*{(1.3)}
\]
Soit $S_\alpha\inserted{=\operatorname{Spec}V_\alpha}$ le normalisé de $S$ dans $K_\alpha$, qui est encore \unsure{str.\textsuperscript{t}} local, \deleted{\unsure{soit}} et soit $s_\alpha$ le pt fermé de $S_\alpha$, $\eta_\alpha$ son pt \unsure{générique}. \inserted{\unsure{Notons}} et
\[
X_\alpha=X\times_S S_\alpha,\quad X_{\alpha0}=X_{\alpha s_\alpha}=X_0\otimes_k k_\alpha
\tag*{(1.4)}
\]
